\documentstyle[%


righttag%


,11pt%


,amssymb%


,russian%


,izvru%


]{amsart}





\newcommand{\tl}[3]{\medskip\noindent{\bf#1}\ #2 \dotfill #3 \par} 


\pagestyle{empty}


% \pagestyle{plain}


\begin{document}


\par


\vskip 2cm


\par


\bigskip





Номер 12 журнала ``Известия вузов. Математика'' за 2001 г. будет


посвящен задачам оптимального управления и математического


программирования.


\bigskip





\bigskip


\par


\centerline{\Large СОДЕРЖАНИЕ}


\bigskip


\bigskip


%%%%%%%%%% Use \newline %%%%%%%%%%%








\tl{Абубакиров Н.Р., Салимов Р.Б., Шабалин П.Л.}


{Внешняя обратная краевая задача\newline при комбинировании двух параметров из


декартовых координат и полярного угла}{3}


\tl{Галиуллин И.А.} 


{Бэровский класс показателей Ляпунова механических 


систем,\newline содержащих параметры}{11}


\tl{Данеев А.В., Русанов В.А.} 


{Геометрический подход к решению некоторых обратных \newline


задач системного анализа}{18}


\tl{Драгилев М.М.}


{О пространствах К\"ете 


с предэквивалентными безусловными базисами}{29}


\tl{Жигалко Ю.П., Соловьев С.И.} 


{Собственные колебания балки с гармоническим\newline осциллятором}{36}


\tl{Климов В.С.} 


{Метод усреднения в задаче о периодических решениях 


квазилинейных\newline параболических уравнений}{39}


\tl{Кытманов А.М., Мысливец М.С.} 


{О $CR$-распределениях, заданных на\newline гиперповерхности}{47}


\tl{Майборода И.Н., Островецкий Л.А.} 


{О неподвижных точках операторов\newline с гетеротонной производной}{53}


\tl{Павлова М.Ф., Шемуранова Е.В.}


{О существовании слабого решения одной задачи\newline ненасыщенной


фильтрационной консолидации}{58}


\tl{Покровский В.Г.} 


{О разрезании многомерного параллелепипеда на параллелепипеды}{69}


\tl{Сабитов К.Б., Акимов А.А.} 


{К теории аналога задачи Неймана для уравнений\newline смешанного типа}{73}








\vfill


\noindent\raisebox{2pt}{\copyright} Казанский госудаpственный унивеpситет


\end{document}


